\begin{abstract}
Low Earth Orbit (LEO) operations now face simultaneous stress from mega-constellation growth and the Solar Cycle 25 activity peak in 2025--2026. Under these conditions, the Simplified General Perturbations-4 (SGP4) propagator suffers severe degradation because static drag terms cannot track storm-time thermospheric expansion. We present a Gated Physics-Informed Neural Network (G-PINN) that learns residual corrections while enforcing a hard physical boundary condition through a differentiable Physics Gate, $\Gamma(t)=\tanh(\lambda t)$. This gate guarantees zero correction at the TLE epoch, eliminating the non-physical Phantom Drift commonly observed in unconstrained ML correctors. On curated Starlink trajectories with fused space-weather features ($K_p$ and $F_{10.7}$), the model improves storm-regime accuracy, reducing mean error from 11.18 km to 9.77 km and compressing worst-case search radius from approximately 28 km to 21 km. The learned corrections also reveal a persistent $+5.5$ m/s radial velocity bias in baseline SGP4 handling, demonstrating diagnostic value beyond forecasting. We integrate these corrections into a Mahalanobis-based Monte Carlo conjunction pipeline and a high-fidelity WebGL visualization stack, enabling resilient and operationally interpretable collision-risk assessment during extreme geomagnetic activity.
\end{abstract}